\section{Introduction}
\label{sec:intro}

Large language models learn the social biases latent in their pretraining data, and in
deployment they do not merely reproduce those biases, they amplify them. The cost is
borne disproportionately by minoritized groups, who are turned into the default
referent whenever a model fills an ambiguous gap with a stereotype. A growing toolkit
of interventions aims to reduce this harm, from data curation and fine-tuning to
mechanistic interpretability and representation engineering. A recurring lesson is that
bias is not a single universal quantity: what counts as a harmful stereotype, and which
intervention will undo it, depends on the social context. Interventions must therefore
be \emph{context-aware}.

We take a first step toward adapting such interventions to the Latin American (LATAM)
context. Doing so was, until recently, blocked by a simpler problem: there was no
Spanish-language data with which to even measure social bias in the region. The new
\textsc{SESGO} benchmark~\citep{robles2025sesgo} closes that gap with
Latin-American-grounded items, and we leverage it to characterize bias in
Spanish-prompted LLMs through behavioral and distributional studies. Our
goal is not a finished mitigation but a map: a preliminary, multi-view picture of where
bias lives in these models that future LATAM-specific interventions can build on.

A LATAM-specific consideration shapes what we measure: \textbf{model size}. Adaptation
in the region happens predominantly through \emph{open-weight} models, both because such
models are the substrate practitioners can actually fine-tune and because open weights
are increasingly seen as a precondition for AI sovereignty~\citep{sovereignty}.
Resource constraints push real-world adopters toward the smallest models, which, as we
find, can also be the most biased, so we study how bias varies with scale across open
Qwen, Llama, Mistral, and Gemma families.

\paragraph{Preliminary findings.} Across these studies, three trends track model scale:
\begin{itemize}[leftmargin=*,topsep=2pt,itemsep=1pt]
  \item \textbf{Smaller models are more biased:} abstention accuracy on ambiguous items rises
    with scale, and the residual error split toward the marginalized group narrows
    (\Cref{sec:bias}).
  \item \textbf{Smaller models answer less consistently:} the smallest model in each family
    flips its answer under cosmetic reformatting, while the largest does not
    (\Cref{sec:stability}).
  \item \textbf{Reasoning commits at a few discrete tokens:} forking the thinking models' chain
    of thought, every model from $0.8$B to $9$B resolves its answer at a small number of decisive
    forking tokens (change-point Bayes factor ${\approx}10^{9}$) rather than accreting it
    gradually (\Cref{sec:uncertainty}).
\end{itemize}
These are preliminary: the bias and stability results span the full sweep, while the
reasoning-dynamics results are single-item case studies, offered as a foundation rather than a
conclusion.

\paragraph{Contributions.} Building on \textsc{SESGO}~\citep{robles2025sesgo}, we contribute:
\begin{itemize}[leftmargin=*,topsep=2pt,itemsep=1pt]
  \item \textbf{An extension of \textsc{SESGO} to model scale,} replicating the benchmark
    across open-weight Qwen, Llama, Mistral, and Gemma models to chart how Spanish-prompted
    bias changes with size.
  \item \textbf{A stability analysis} of how consistently models answer under cosmetic
    reformatting.
  \item \textbf{A chain-of-thought dynamical study of bias commitment,} tracking how and when
    an open-weight model decides its answer to an ambiguous item.
\end{itemize}

\paragraph{Setup.} We evaluate the open-weight sweep in \Cref{tab:models} on the full Spanish
\textsc{SESGO} prompt set, reading each item with a greedy decode parsed to a committed choice
(and a three-option teacher-forced distribution over the roles $\{t,o,u\}$ for the answer-token
probability); all rates carry Wilson $95\%$ CIs. The three sections that follow study how model
\emph{scale} shapes, in turn, answer stability (\Cref{sec:stability}), \textsc{SESGO} bias
(\Cref{sec:bias}), and reasoning uncertainty (\Cref{sec:uncertainty}).

\begin{table}[ht]
  \centering \small
  \caption{The open-weight model sweep.}
  \label{tab:models}
  \begin{tabular}{@{}lll@{}}
    \toprule
    \textbf{Family} & \textbf{Sizes} & \textbf{Variants} \\
    \midrule
    \texttt{Qwen3.5}       & $0.8$, $2$, $4$, $9$, $27$B & non-thinking $+$ thinking \\
    \texttt{Gemma-4}       & E2B, E4B                    & instruct \\
    \texttt{Ministral-3}   & $3$, $8$B                   & Instruct $+$ Reasoning \\
    \texttt{Llama-3.1/3.2} & $1$, $3$, $8$B              & instruct \\
    \bottomrule
  \end{tabular}
\end{table}
